 7(1-y)-6\min(\alpha,1-y)-6(1+\alpha-2y)_+
       +(1+4\alpha-2y)_+,             &1/2<y<1,\\
 0,                                  &y>1.
 \end{cases}
\end{equation}
Endpoint values may be assigned arbitrarily for integration. The rank
correction is
\begin{equation}\label{eq:douter}
 d(y)=\begin{cases}
 \bigl(1+4\alpha-3y-(1+\alpha-3y)_+\bigr)_+,&1/3<y<1/2,\\
 0,&\text{otherwise}.
 \end{cases}
\end{equation}
This is exactly source (5.9), with its \emph{multiplicative} indicator
written piecewise. In particular
$\int_{1/3}^{1/2}d(y)\,dy=9/640$.

Resolving \eqref{eq:gammaout} and using the fact that all integer endpoint
corrections are bounded gives, uniformly for outer primes $p\le K$,
\begin{align}
 -\gamma_p^{\mathrm{out}}
   &=K\bigl(R_0(p/K)-d(p/K)\bigr)+O(1),\label{eq:outergammalimit}\\
 -v_p(S_K)&=K\left(-2\lambda\floor{1/y}
                   +\sum_{j=1}^5(2\lambda-jy)_+\right)+O(1).
 \label{eq:outerscalarlimit}
\end{align}
The scalar formula holds throughout $1/3<y\le2\lambda$; $N<p$ and
$p^2>5K$ in this range. To verify uniformity directly, use
$u=K(\alpha+1-(\floor{1/y}+1)y)_++O(1)$ and
$t_p=K\min(\alpha,1-\floor{1/y}y)+u$, then the Lipschitz property of
minima and positive parts. The rank minimum in \eqref{eq:gammaout}
is exactly what produces the subtraction of $d$.

Let
\begin{equation}\label{eq:Iout}
 I_{\mathrm{out}}=\int_{1/3}^{2\lambda}
 \left(R_0(y)-d(y)-2\lambda\floor{1/y}
                   +\sum_{j=1}^5(2\lambda-jy)_+\right)\,dy
 =\frac{127751}{96000}.
\end{equation}
The exact piecewise evaluation is in Appendix~\ref{app:arithmetic}.
For $1<y<2\lambda$, both residue terms vanish but the factorial scalar
still contributes; this interval must not be omitted.

\begin{lemma}[Weighted prime-sum consequence of the PNT]
\label{lem:primesuminterface}
Let $0<a<b<\infty$, and let $f$ be bounded and piecewise continuous on
$[a,b]$. Then
\[
 \frac1{K^2}\sum_{aK<p\le bK}Kf(p/K)\log p
       \longrightarrow\int_a^b f(y)\,dy.
\]
Consequently, for bounded piecewise continuous $g$ on $[3,M]$,
\[
 \frac1{K^2}\sum_{K/M<p\le K/3}p\,g(K/p)\log p
       \longrightarrow\int_3^M g(x)x^{-3}\,dx.
\]
\gid{Zeta5.weighted_prime_sum}
\end{lemma}
\begin{proof}
Equation~\eqref{eq:pnt} implies
$(\vartheta(vK)-\vartheta(uK))/K\to v-u$ for fixed $0<u<v$.
This proves the first statement for interval step functions. Upper and
lower step approximations on a finite partition prove it for bounded
Riemann-integrable functions, including bounded piecewise continuous
ones. Bounds on the total measure follow from
$\vartheta(bK)/K=O(1)$. A value at any one endpoint contributes at most
$O(\log K/K)$, so endpoint conventions do not matter. For the second
statement use $f(y)=y g(1/y)$ on $[1/M,1/3]$ and the substitution $x=1/y$.
\end{proof}

\begin{proposition}[Normalization growth; source Proposition 5.2]
\label{prop:primegrowth}
For every fixed $M\ge40$,
\begin{equation}\label{eq:primegrowth}
 \limsup_{\substack{K\to\infty\\40\mid K}}
 K^{-2}\log m_{K,M}
 \le I_{\mathrm{out}}+\frac{6\lambda}{M}
                         +\int_3^M R(x)x^{-3}\,dx.
\end{equation}
\gid{Zeta5.normalization_growth}
\end{proposition}
\begin{proof}
Taking logarithms of \eqref{eq:mdef} gives
$\log m_{K,M}=-\sum_{p\le2h}L_p(K,M)\log p$, with signs unchanged.
For $p\le\sqrt{5K}$ the bound \eqref{eq:Fsmall} contributes at most
$O(K^{3/2}\log K)=o(K^2)$; the factors $v_p(24)$ involve only $2,3$.
For $\sqrt{5K}<p\le K/M$ the integer logarithm is one, and the PNT gives
$6\lambda K^2/M+o(K^2)$.

In the inner range, Lemma~\ref{lem:innerlimit} gives
$-L_p=pR(K/p)+O_M(1)$. Its total error is bounded by
$O_M(\vartheta(K))=O_M(K)=o(K^2)$. Apply
Lemma~\ref{lem:primesuminterface}. The outer estimates
\eqref{eq:outergammalimit}--\eqref{eq:outerscalarlimit} give
\eqref{eq:Iout} by the same lemma. All constants depending on $M$ are
used with $M$ fixed before $K\to\infty$.
\end{proof}

\subsection{The periodic tail and the two fixed cutoffs}
For $x\ge3$ write
\[
 f=\{x\},\quad g=\{\alpha x\},\quad \tau=\{2Hx\},\quad
 \sigma=\{2x\},\quad\eta=\{2\lambda x\},
\]
and let $A(z)=\ell(x,z)-2x$, $B(z)=\ell(\alpha x,z)-2\alpha x$.
Expansion of \eqref{eq:Gamma}--\eqref{eq:Rlimit} gives
\begin{equation}\label{eq:RFQ}
 R(x)=xF(x)+Q(x),\qquad F(x)=4\lambda+2\lambda f-12\lambda g.
\end{equation}
\begin{equation}\label{eq:Qtail}
\begin{split}
 Q(x)={}&\frac{\tau(\tau-\sigma)-(\tau-\sigma)_++\eta(1-\eta)}2\\
 &+9\int_0^{1/2}B(z)^2\,dz-3\int_0^{1/2}A(z)B(z)\,dz.
\end{split}
\end{equation}
Both $A$ and $B$ have integral zero on $(0,1/2)$, and each squared
integral is at most $1/8$. The mixed integral has absolute value at most
$1/8$ by Cauchy--Schwarz. The numerator before the division by two in
\eqref{eq:Qtail} is between $-1/4$ and $1/4$: if $\tau\le\sigma$, its
first two terms equal $\tau(\tau-\sigma)\in[-1/4,0]$; otherwise they
are $-(\tau-\sigma)(1-\tau)\in[-1/4,0]$, and
$\eta(1-\eta)\in[0,1/4]$. Hence
\begin{equation}\label{eq:Qtailbounds}
 -\frac12\le Q(x)\le\frac{13}{8}.
\end{equation}
An explicit exact verification of the expansion on every finite branch
used below is supplied with Appendix~\ref{app:arithmetic}.

Define continuous periodic functions
\begin{align*}
 \mathcal P(x)&=74g(1-g)-\lambda f(1-f),&
 \overline{\mathcal P}&=\frac{2923}{240},\\
 G_0(v)&=v(1-v)(2v-1)/6,&
 \mathcal C(x)&=(74/\alpha)G_0(g)-\lambda G_0(f).
\end{align*}
Off their locally finite breakpoints,
$\mathcal P'=F+\lambda$ and
$\mathcal C'=\mathcal P-\overline{\mathcal P}$. They are continuous at
those breakpoints, so integration by parts over subintervals telescopes
without jump terms. Since
$\max_{[0,1]}|G_0|=1/(36\sqrt3)$,
\[
 |\mathcal C(x)|\le\frac{118511}{4320\sqrt3}<16.
\]
Two integrations by parts, first on a finite interval and then at
infinity, give
\begin{equation}\label{eq:tailidentity}
\begin{split}
 \int_T^\infty\frac{R(x)}{x^3}\,dx={}&-\frac\lambda T
 -\frac{\mathcal P(T)}{T^2}+\frac{\overline{\mathcal P}}{T^2}
 -\frac{2\mathcal C(T)}{T^3}\\
 &+6\int_T^\infty\frac{\mathcal C(x)}{x^4}\,dx
 +\int_T^\infty\frac{Q(x)}{x^3}\,dx.
\end{split}
\end{equation}
The boundary limits and absolute convergence follow from boundedness of
$F,Q,\mathcal P,\mathcal C$.

At $T=20$, $\mathcal P(T)=37/2$ and $\mathcal C(T)=0$, so
\begin{equation}\label{eq:tail20}
 \int_{20}^\infty R(x)x^{-3}\,dx\le-\frac{2689}{48000}.
\end{equation}
If $40\mid M$, then $\mathcal P(M)=\mathcal C(M)=0$, and
\begin{equation}\label{eq:tailM}
 \int_M^\infty R(x)x^{-3}\,dx
 \ge-\frac\lambda M+\frac{2923/240-1/4}{M^2}-\frac{32}{M^3}.
\end{equation}
The finite integral is
\begin{equation}\label{eq:innerexact}
 \int_3^{20}R(x)x^{-3}\,dx
 =\frac{322437603634266857629}{7535670527041937280000}.
\end{equation}
Put
\begin{align}
 A_*&=\frac{9928298118277006344769}{7535670527041937280000},
 \label{eq:Astar}\\
 A_M&=A_*+\frac{7\lambda}{M}
              -\frac{2923/240-1/4}{M^2}+\frac{32}{M^3}.
 \label{eq:AM}
\end{align}
Combining Proposition~\ref{prop:primegrowth} with
\eqref{eq:tail20}--\eqref{eq:innerexact} proves
\begin{equation}\label{eq:AMbound}
 \limsup_{\substack{K\to\infty\\40\mid K}}K^{-2}\log m_{K,M}\le A_M
 \qquad(M\in40\Z_{>0}).
\end{equation}
\gid{Zeta5.normalization_quadratic_bound}

\section{The real determinant estimate}
\label{sec:real}
Define the external field on $[0,\infty)$ by
\begin{equation}\label{eq:Vdef}
 V(t)=2\pi\sqrt t+\int_0^1\log(t+u^2)\,du
                         -6\int_0^\alpha\log(t+u^2)\,du.
\end{equation}
The integrals at $t=0$ are interpreted as their convergent improper
integrals. For a finite measure $\rho$, write
\[
 U^\rho(t)=\int\log|t-u|\,d\rho(u),\qquad
 I(\rho)=\iint\log|t-u|\,d\rho(t)\,d\rho(u),
\]
whenever the integrals are absolutely convergent. The same notation is
used on $\C$.

\begin{proposition}[Explicit comparison measure; source Lemma 6.1]
\label{prop:comparison}
There is a positive measure $\rho$ of mass $\lambda$, supported in
$(0,2)$ and equal to the sum of the sixteen rationally specified arcsine
measures in Appendix~\ref{app:measure}, such that
\begin{equation}\label{eq:potentialbound}
 2U^\rho(t)-V(t)\le M_0:=-\frac{1329}{200}\qquad(t\ge0).
\end{equation}
Every supporting interval has length greater than $1/225$. With
\begin{align}
 C_*&=-2\lambda+12\alpha\lambda(1-\log\alpha)
                  +3\lambda^2-2\lambda^2\log(2\lambda),\label{eq:Cstar}\\
 U&=-\frac{2733991}{2000000},\label{eq:Uconstant}
\end{align}
one has
\begin{equation}\label{eq:energyconstant}
 \lambda M_0-I(\rho)+C_*<U.
\end{equation}
\gid{Zeta5.comparison_measure_certificate}
\end{proposition}
\begin{proof}
See Appendix~\ref{app:measure}. The proof uses a finite rational
partition with outward bounds for elementary functions. The independent
checker supplied with this edition verifies all its cells and constants.
This proposition uses only a comparison measure, not a solution or an
existence theorem for the equilibrium-measure problem.
\end{proof}

\subsection{Zero-mass energy with explicit integrability}
\begin{lemma}[Zero-mass logarithmic energy; source Lemma 6.2]
\label{lem:zeroenergy}
Let $\nu$ be a finite real signed Borel measure on $\C$ with compact
support, $\nu(\C)=0$, and
\[
 \iint|\log|z-w||\,d|\nu|(z)\,d|\nu|(w)<\infty.
\]
Here the absolute logarithmic kernel is assigned the value $+\infty$ on
the diagonal $z=w$. In the signed real integral, use any representative
on the resulting null diagonal. Then $I(\nu)\le0$.
\gid{Zeta5.Analysis.zero_mass_logarithmic_energy}
\end{lemma}
\begin{proof}
For $s>0$, completing the square gives
\[
 e^{-s|z-w|^2}=\frac{4s}{\pi}\int_\C
          e^{-2s|z-u|^2}e^{-2s|w-u|^2}\,dA(u),
\]
where $A$ is planar Lebesgue measure. Thus
\[
 J(s):=\iint e^{-s|z-w|^2}\,d\nu(z)\,d\nu(w)
 =\frac{4s}{\pi}\int_\C
       \left(\int e^{-2s|z-u|^2}\,d\nu(z)\right)^2dA(u)\ge0.
\]
Fubini is justified by first integrating the nonnegative kernel against
$|\nu|\otimes|\nu|\otimes A$; the completed-square identity bounds that
integral by a constant times $|\nu|(\C)^2$.

For $0<a<b$ and $r>0$, put
\[
 L_{a,b}(r)=\frac12\int_a^b\frac{e^{-s}-e^{-sr^2}}s\,ds.
\]
The identity
$\frac12\int_0^\infty(e^{-s}-e^{-sr^2})s^{-1}\,ds=\log r$
follows by writing the numerator as $\int_1^{r^2}s e^{-sv}\,dv$ and
integrating. The integrand has a fixed sign for fixed $r$, so
$|L_{a,b}(r)|\le|\log r|$ and $L_{a,b}(r)\to\log r$ as
$a\downarrow0$, $b\uparrow\infty$. Since $\nu$ has mass zero,
\[
 \iint L_{a,b}(|z-w|)\,d\nu(z)\,d\nu(w)
       =-\frac12\int_a^bJ(s)s^{-1}\,ds\le0.
\]
The absolute logarithmic-integrability hypothesis implies that the
diagonal has zero $|\nu|\otimes|\nu|$ measure. Dominated convergence
therefore applies off that null set and proves the conclusion.
\end{proof}

\begin{lemma}[Circle regularization and arcsine error]
\label{lem:regularization}
For $t_1,\ldots,t_h\ge0$ and $\varepsilon>0$, let $\omega_i$ be
normalized arclength on the circle of radius $\varepsilon$ centered at
$t_i$, and set $\sigma=K^{-1}\sum_{i=1}^h\omega_i$. Then
\begin{equation}\label{eq:configuration}
\begin{split}
 2\sum_{i<j}\log|t_i-t_j|-K\sum_iV(t_i)
 \le{}&(\lambda M_0-I(\rho))K^2\\
 &+120\lambda K^2\sqrt\varepsilon-h\log\varepsilon.
\end{split}
\end{equation}
Coincident points are allowed by assigning $-\infty$ to the left side.
\gid{Zeta5.Analysis.regularized_configuration_bound}
\end{lemma}
\begin{proof}
The circle-average identity is
\[
 \int\log|z-u|\,d\omega_i(z)=\log\max(|t_i-u|,\varepsilon).
